\section{Discussion}
\label{sec:discussion}

\subsection{When does quantization pay off?}
Our measurements reduce to five rules for deploying a quantized LLM on a GPU like the RTX 5090.

\textbf{(1) Budget memory with the real ratio, not $16/b$.} Every format saves memory, but the
16-bit embedding, LM head and scales cap 4-bit formats at \WtNFSeven--\WtAWQSeven{} for a 7B
model, and the KV cache and activations are not quantized at all. A capacity plan that assumes
$4\times$ will run out of memory.

\textbf{(2) For interactive, low-concurrency serving, use 4-bit weight-only formats with fused
kernels.} At one request per GPU, AWQ-INT4 with Marlin decodes \AWQspeedup{} faster than BF16
and uses \EnRatioAWQone{} less energy per token.

\textbf{(3) For high-throughput serving, use W8A8 FP8.} Weight-only formats lose their advantage
once the batch makes decode compute-bound ($B\gtrsim\AWQcross$ for AWQ at 512-token prompts) and
never speed up prefill. FP8 is the only format that keeps winning, with \ThrFPratio{} the BF16
throughput at $B=256$, because it also doubles the tensor-core rate.

\textbf{(4) Quantize for memory, not speed, when only bitsandbytes is available.} LLM.int8() is
\INTeightslowdown{} slower than BF16 at batch size one and never catches up, and NF4 is faster only
for small batches of a model large enough to be memory-bound. Both remain useful to fit a model
that otherwise would not fit.

\textbf{(5) \GuideHead} \QualityGuide

\noindent Across all five, the bit-width alone predicts very little; the kernel that executes the
format and the regime the workload puts the GPU in decide the outcome.

\textbf{Limitations.} We test one GPU and dense Qwen models. The mechanism is architectural and Qwen3
reproduces it, but crossover batch sizes depend on each GPU's ratio of compute to bandwidth. The GPU
also drives a display, so we could not lock its clocks (repeated runs still agree closely). We do not
study 4-bit activations, KV-cache quantization or mixture-of-experts models.
\label{sec:limits}


\section{Conclusion}
\label{sec:conclusion}
We measured how much memory LLM quantization saves and when it makes inference faster, using
eight configurations of Qwen2.5 models on an RTX 5090. Memory savings are real but smaller than
advertised: \WtNFSeven{} instead of $4\times$ for 4-bit weights, and as little as \MemBigNF{} end
to end once the KV cache grows. Speed depends on the kernel and the regime rather than on the
bit-width: the same 8-bit weights decode \GPTQeightspeedup{} faster than BF16 with Marlin and
\INTeightslowdown{} slower with bitsandbytes, weight-only gains disappear in the compute-bound regime, and FP8 W8A8 is the only
format that is faster everywhere. \QualityConclusion{} Smaller is not always faster---but with the
right kernel, for the right regime, it is both smaller and faster.

\section*{Division of work}
Following the proposal, Sijia Zhong built the baseline harness and the memory study (Q1), Zheng Zheng
ran the controlled precision/batch/length sweeps and the kernel and profiling analysis (Q2), and
Haofei Sun ran the ShareGPT and quality evaluations (Q3, Q4). All three analyzed the results and
wrote the report and slides.
